%% =====================================================================
%%  Unit 17 — Where were you?  (Parler d'hier)
%%  Pistes : 17.01–17.04 dialogues · 17.05–17.06 vocabulaire · 17.07–17.08 grammaire
%%           17.09 prononciation · 17.10–17.13 exercices (textes en annexe)
%% =====================================================================

\newcommand{\otIllusUnXVII}{\otIllusDuo{yesterday/104/86/15/0.10,ago/190/92/20/0.09,last week/122/14/13/0.08,
  was/150/8/11/0.08,born/196/22/12/0.07,were/100/50/12/0.07}{jack}{Where were you?}{lea}{I was in Times Square!}}

% Frise chronologique de New York (exercice 5)
\newcommand{\otFriseNY}{%
  \begin{tikzpicture}[x=1mm,y=1mm]
    \draw[otNavy,line width=1.6pt,-{Stealth[length=3mm]}] (0,0) -- (164,0);
    \foreach \x/\an/\txt/\c/\h in {15/1625/{New Amsterdam:\\a Dutch town}/otTeal/1,
      41/1664/{The English arrive:\\New Amsterdam\\becomes New York}/otRedDark/-1,
      67/1883/{The Brooklyn\\Bridge}/otOrangeDark/1,
      93/1886/{The Statue of Liberty,\\a gift from France}/otGreenDark/-1,
      119/1892/{Ellis Island:\\the door to America}/otViolet/1,
      145/1931/{The Empire State\\Building}/otNavy/-1}{%
      \fill[\c] (\x,0) circle (2.2);
      \fill[white] (\x,0) circle (0.9);
      \node[font=\otBold\fontsize{10.5}{12}\selectfont,text=\c] at (\x,5*\h) {\an};
      \ifnum\h=1 \def\otAncre{south}\else\def\otAncre{north}\fi
      \node[font=\fontsize{8.4}{10}\selectfont,align=center,text=otInk,anchor=\otAncre]
        at (\x,8.6*\h) {\txt};}
  \end{tikzpicture}}

\unite[Where were\\you?]{17}{Where were you?}{Parler d'hier}{otIllusUnXVII}

\situation{Épisode 17}{Lundi 5 janvier. Tout le monde est de retour au studio après les fêtes, et
chacun raconte où il était pour le Nouvel An. Le dimanche suivant, Mrs.\ Lopez montre à Léa de vieilles
photos du quartier.}

\begin{objectifs}
\item dire où vous étiez et comment c'était, avec \en{was} et \en{were} ;
\item poser des questions sur le passé et y répondre ;
\item décrire un lieu autrefois avec \en{there was}, \en{there were} ;
\item situer un moment dans le passé : \en{yesterday}, \en{last week}, \en{two years ago} ;
\item dire où et quand vous êtes né ;
\item prononcer les formes faibles de \en{was} et \en{were}.
\end{objectifs}

\auprogramme
  {Dire où l'on était, décrire le passé}
  {\en{was}, \en{were}\par\vspace{0.6mm}\en{there was}, \en{there were} ; \en{ago}}
  {Hier et autrefois : 24 expressions}
  {Les formes faibles de \en{was} et \en{were}}
  {Une brève histoire de New York}

\semaine{casque/otNavy/Dialogues/Écouter et lire les deux scènes,
  bulle/otGreenDark/Vocabulaire/Apprendre les 24 expressions,
  blocs/otViolet/Grammaire/was et were et ago,
  sons/otOrangeDark/Sons et culture/Formes faibles et histoire,
  crayon/otNavy/Exercices/12 exercices corrigés,
  coche/otGreenDark/Bilan/L'essentiel et le test express}

\clearpage
%% ---------------------------------------------------------------------
\rubrique{ecouter}{Souvenirs, souvenirs}

\begin{dialogue}{1}{How was London?}{Lundi 5 janvier, 9 heures. Emma rentre de Londres ; Jack et Léa
  l'attendent avec des cafés.}{17.01}{17.02}
\dit{jack}{Welcome back, Emma! \hl{How was} London?\note{1}}{Bon retour, Emma ! C'était comment, Londres ?}
\dit{emma}{\hl{It was} great! My family \hl{was} so happy. But the weather \hl{was} terrible!}{C'était super ! Ma famille était si contente. Mais il faisait un temps affreux !}
\dit{lea}{And \hl{where were you} on New Year's Eve?\note{2}}{Et où étais-tu le soir du Nouvel An ?}
\dit{emma}{I \hl{was} at a party with my brother. And \hl{where were} you?}{J'étais à une fête avec mon frère. Et toi, où étais-tu ?}
\dit{lea}{I \hl{was} in Times Square with my parents! \hl{There were} about a million people!\note{3}}{J'étais à Times Square avec mes parents ! Il y avait environ un million de personnes !}
\dit{emma}{Great! \hl{Were} your parents happy?}{Génial ! Tes parents étaient contents ?}
\dit{lea}{\hl{Yes, they were}. But they \hl{were} very cold!}{Oui. Mais ils avaient très froid !}
\dit{emma}{And you, Jack? \hl{Were you} in Times Square too?}{Et toi, Jack ? Tu étais à Times Square, toi aussi ?}
\dit{jack}{\hl{No, I wasn't}. I \hl{was} at home… in bed, at ten!\note{4}}{Non. J'étais chez moi… au lit, à dix heures !}
\end{dialogue}

\begin{comprendre}
\pointnote{1} \en{\textbf{How was London?}} \en{was} est le prétérit (le passé) de \en{is} : « c'était
  comment ? ». Avec \en{I}, \en{he}, \en{she}, \en{it}, on emploie \en{was}.
\pointnote{2} \en{\textbf{Where were you?}} \en{were} est le passé de \en{are} : on l'emploie avec
  \en{you}, \en{we}, \en{they}. « Où étais-tu ? »
\pointnote{3} \en{\textbf{There were…}} Le passé de \en{there are} : « il y avait » (au singulier :
  \en{there was}).
\pointnote{4} \en{\textbf{No, I wasn't.}} La négation : \en{wasn't} (= \en{was not}), \en{weren't}
  (= \en{were not}). La réponse courte reprend \en{was} ou \en{were}.
\end{comprendre}

\begin{dialogue}{2}{A long time ago}{Dimanche après-midi, chez Mrs.\ Lopez. Sur la table, une vieille
  boîte pleine de photos en noir et blanc.}{17.03}{17.04}
\dit{lopez}{Look, Léa! This \hl{was} our street \hl{fifty years ago}.\note{5}}{Regardez, Léa ! C'était notre rue, il y a cinquante ans.}
\dit{lea}{Wow! \hl{There wasn't} a coffee shop!}{Ouah ! Il n'y avait pas de café !}
\dit{lopez}{No, but \hl{there was} a small bakery. And \hl{there were} kids everywhere!}{Non, mais il y avait une petite boulangerie. Et il y avait des enfants partout !}
\dit{lea}{\hl{Where were you born}, Mrs.\ Lopez?\note{6}}{Où êtes-vous née, Mrs.\ Lopez ?}
\dit{lopez}{\hl{I was born} here, in Brooklyn, in nineteen sixty-two.}{Je suis née ici, à Brooklyn, en 1962.}
\dit{lea}{And your parents?}{Et vos parents ?}
\dit{lopez}{They \hl{were} from Puerto Rico. They \hl{were} very young, and very poor.}{Ils venaient de Porto Rico. Ils étaient très jeunes, et très pauvres.}
\dit{lea}{\hl{Were you} happy here?}{Vous étiez heureuse, ici ?}
\dit{lopez}{\hl{Yes, I was}. Life \hl{wasn't} easy, but we \hl{were} happy.}{Oui. La vie n'était pas facile, mais nous étions heureux.}
\end{dialogue}

\begin{comprendre}
\pointnote{5} \en{\textbf{fifty years ago}} \en{ago} se place \textbf{après} la durée : \en{fifty years
  ago} = « il y a cinquante ans » ; \en{two days ago} = « il y a deux jours ».
\pointnote{6} \en{\textbf{Where were you born?}} « Où êtes-vous né(e) ? ». En anglais, « naître » se dit
  \en{be born}, au passé : \en{I was born in Lyon}, \en{I was born in 1998}.
\end{comprendre}

%% ---------------------------------------------------------------------
\rubrique{mots}{Les mots de l'unité}

Voici les 24 expressions de l'unité. Notez que \en{last} (« dernier ») et \en{ago} (« il y a »)
s'emploient sans \en{the} ni préposition : \en{last week}, \en{two days ago}.

\begin{vocabulaire}{Hier et autrefois}{17.05}
\categorie{Quand ?}
\mot{yesterday}{ˈjestɚdeɪ}{hier}
\mot{yesterday morning}{ˌjestɚdeɪ ˈmɔːrnɪŋ}{hier matin}
\mot{last night}{ˌlæst ˈnaɪt}{hier soir, la nuit dernière}
\mot{last week}{ˌlæst ˈwiːk}{la semaine dernière}
\mot{last year}{ˌlæst ˈjɪr}{l'année dernière}
\mot{two days ago}{ˌtuː ˈdeɪz əˈɡoʊ}{il y a deux jours}
\mot{a long time ago}{ə ˌlɑːŋ ˈtaɪm əˈɡoʊ}{il y a longtemps}
\mot[in nineteen ninety]{in 1990}{ɪn ˌnaɪntiːn ˈnaɪnti}{en 1990}
\categorie{Parler du passé}
\mot{Where were you?}{ˌwer wɚ ˈjuː}{Où étais-tu ?}
\mot{I was at home.}{aɪ wəz ət ˈhoʊm}{J'étais à la maison.}
\mot{How was it?}{ˌhaʊ ˈwʌz ɪt}{C'était comment ?}
\mot{Where were you born?}{ˌwer wɚ juː ˈbɔːrn}{Où es-tu né(e) ?}
\mot{I was born in France.}{aɪ wəz ˌbɔːrn ɪn ˈfræns}{Je suis né(e) en France.}
\end{vocabulaire}

\begin{vocabulaire}{La ville d'autrefois}{17.06}
\categorie{L'histoire}
\mot{history}{ˈhɪstəri}{l'histoire (le passé)}
\mot{a century}{ə ˈsentʃəri}{un siècle}
\mot{an immigrant}{ən ˈɪmɪɡrənt}{un immigrant}
\mot{a building}{ə ˈbɪldɪŋ}{un immeuble, un bâtiment}
\mot{a statue}{ə ˈstætʃuː}{une statue}
\mot{a gift}{ə ˈɡɪft}{un cadeau}
\categorie{Pour décrire}
\mot{famous}{ˈfeɪməs}{célèbre}
\mot{tall}{tɑːl}{haut, grand}
\mot{poor}{pʊr}{pauvre}
\mot{rich}{rɪtʃ}{riche}
\mot{crowded}{ˈkraʊdɪd}{bondé, plein de monde}
\end{vocabulaire}

%% ---------------------------------------------------------------------
\rubrique[was et were]{grammaire}{\en{was} et \en{were} : le passé de \en{be}}

\otIntertitre[otViolet]{1 · \en{was} et \en{were}}
\begin{center}\otTableau
\begin{tabular}{@{}l@{\hspace{8mm}}>{\enOn}l<{\enOff}@{\hspace{8mm}}l@{}}
\otTH{Forme} & \otTH{Exemple} & \otTH{Français}\\\midrule
affirmation      & I / he / she / it \textbf{was} at home. · You / we / they \textbf{were} at home. & j'étais…\\
négation         & I \textbf{wasn't} late. · They \textbf{weren't} late.                              & je n'étais pas…\\
question         & \textbf{Was} she happy? · \textbf{Were} you at work?                               & était-elle… ?\\
réponses courtes & Yes, she \textbf{was}. · No, we \textbf{weren't}.                                  & oui / non\\\bottomrule
\end{tabular}\end{center}

\begin{regle}[Un seul verbe, deux formes]
\en{was} avec \en{I}, \en{he}, \en{she}, \en{it} (et un nom au singulier) ; \en{were} avec \en{you},
\en{we}, \en{they} (et un nom au pluriel). Selon le sens, \en{was} se traduit par « étais », « était »,
« a été », « fut » : \en{It was great!} = « C'était super ! ».
\end{regle}

\blocaudio{17.07}{Écoutez et répétez}{Six phrases avec \en{was} et \en{were}.}

\otIntertitre[otViolet]{2 · \en{there was}, \en{there were} et \en{ago}}
\begin{center}\otTableau
\begin{tabular}{@{}l@{\hspace{8mm}}>{\enOn}l<{\enOff}@{\hspace{8mm}}l@{}}
\otTH{Présent} & \otTH{Passé} & \otTH{Français}\\\midrule
\en{there is} a bakery    & \textbf{there was} a bakery             & il y avait une boulangerie\\
\en{there are} two cinemas & \textbf{there were} two cinemas         & il y avait deux cinémas\\
\en{there isn't} a café   & \textbf{there wasn't} a café            & il n'y avait pas de café\\
\en{Is there…?}           & \textbf{Was there} a park? · \textbf{Were there} any shops? & y avait-il… ?\\\bottomrule
\end{tabular}\end{center}

\begin{regle}[\en{ago} : « il y a » (dans le passé)]
\en{ago} se place \textbf{après} la durée : \en{ten minutes ago}, \en{three years ago}, \en{a long time
ago}. Ne le confondez pas avec \en{there is} (« il y a » = « il existe »).
\end{regle}

\blocaudio{17.08}{Écoutez et répétez}{Une question au passé, puis sa réponse courte.}

\begin{attention}[Les pièges des francophones]
\fauxjuste{I were at home.}{I was at home.}
\fauxjuste{They was happy.}{They were happy.}
\fauxjuste{Did you was at the party?}{Were you at the party?}
\fauxjuste{I am born in 1998.}{I was born in 1998.}
\fauxjuste{There is two years.}{Two years ago.}
\fauxjuste{the last week}{last week}
\end{attention}

%% ---------------------------------------------------------------------
\rubrique[Formes faibles de was et were]{sons}{Les formes faibles de \en{was} et \en{were}}

Dans une phrase, \en{was} et \en{were} sont presque toujours \textbf{faibles} : \en{I was at home} se dit
\phon{aɪ wəz ət ˈhoʊm}, \en{They were late}, \phon{ðeɪ wɚ ˈleɪt}. Ils redeviennent \textbf{forts} en fin de
phrase, dans une réponse courte (\en{Yes, I was} \phon{wʌz}) et à la forme négative (\en{wasn't}
\phon{ˈwʌzənt}, \en{weren't} \phon{wɝːnt}).

\begin{phrases}{17.09}{Écoutez et répétez. Les mots en gras portent l'accent ; \en{was} et \en{were},
eux, restent faibles, sauf en fin de phrase.}
\phrase{I was at home.}{\phon{wəz} · \textbf{home}}
\phrase{It was great!}{\phon{wəz} · \textbf{great}}
\phrase{They were late.}{\phon{wɚ} · \textbf{late}}
\phrase{Where were you?}{\phon{wɚ} · \textbf{where}, \textbf{you}}
\phrase{Yes, I was.}{fin de phrase, forme forte : \phon{wʌz}}
\phrase{No, they weren't.}{négation, forme forte : \phon{wɝːnt}}
\end{phrases}

%% ---------------------------------------------------------------------
\rubrique{culture}{Une brève histoire de New York}

\begin{culture}[De la Nouvelle-Amsterdam à New York]
Vers 1625, des colons hollandais fondent la Nouvelle-Amsterdam à la pointe sud de Manhattan. Selon la
légende, l'île aurait été achetée aux Amérindiens Lenapes pour 60 florins de marchandises. En 1664, les
Anglais prennent la ville et la rebaptisent en l'honneur du duc d'York : New York. Les Hollandais ont
laissé des noms : \en{Brooklyn} (Breukelen), \en{Harlem} (Haarlem), \en{Wall Street} (le mur de la
ville).
\end{culture}

\begin{paire}
\begin{culture}[Ellis Island]
De 1892 à 1954, environ douze millions d'immigrants sont passés par Ellis Island, au pied de la statue
de la Liberté. Aujourd'hui, près de 40\,\% des Américains auraient un ancêtre arrivé par cette porte.
\end{culture}
\begin{culture}[Les cinq \en{boroughs}]
Depuis 1898, New York réunit cinq arrondissements, les \en{boroughs} : Manhattan, Brooklyn, Queens, le
Bronx et Staten Island. Brooklyn compte à lui seul plus de 2,5 millions d'habitants.
\end{culture}
\end{paire}

%% ---------------------------------------------------------------------
\rubrique{pratique}{Exercices}

\exercice[ecrire,niveau=1]{\en{was} ou \en{were} ?}{Complétez avec \en{was} ou \en{were}.}
\begin{anglais}
\begin{questions}[2]
\item I \trou[14mm]{was} at home last night.
\item You \trou[14mm]{were} late yesterday.
\item Emma \trou[14mm]{was} in London.
\item We \trou[14mm]{were} very happy.
\item The party \trou[14mm]{was} great!
\item Léa's parents \trou[14mm]{were} cold.
\item It \trou[14mm]{was} sunny on Sunday.
\item Jack and Kate \trou[14mm]{were} at the park.
\end{questions}
\end{anglais}

\exercice[lire,ecrire,niveau=1]{Hier, à 20 heures}{Où était chacun hier soir ? Écrivez une phrase avec
\en{He was…} ou \en{She was…}.}
\begin{tcbraster}[raster columns=3,raster equal height=rows,raster column skip=3mm,raster row skip=3mm,raster before skip=1mm]
\otCarteAction{lea}{casque}{otGreenDark}{at a concert}{She was at a concert.}
\otCarteAction{jack}{lune}{otNavy}{in bed}{He was in bed.}
\otCarteAction{emma}{globe}{otViolet}{at the airport}{She was at the airport.}
\otCarteAction{sam}{objtasse}{otTeal}{at the café}{He was at the café.}
\otCarteAction{carter}{objtelephone}{otOrangeDark}{at the office}{He was at the office.}
\otCarteAction{lopez}{objsac}{otRed}{at the supermarket}{She was at the supermarket.}
\end{tcbraster}

\exercice[ecrire,niveau=2]{Questions et réponses courtes}{Écrivez la réponse courte qui correspond au
symbole.}
\begin{anglais}
\begin{questions}
\item Were you at home yesterday? \otCheck \hfill \trou[50mm]{Yes, I was.}
\item Was Jack in Times Square? \otCroix \hfill \trou[50mm]{No, he wasn't.}
\item Were Léa's parents in New York? \otCheck \hfill \trou[50mm]{Yes, they were.}
\item Was the weather good in London? \otCroix \hfill \trou[50mm]{No, it wasn't.}
\item Was Emma at a party? \otCheck \hfill \trou[50mm]{Yes, she was.}
\item Were you and Sam late? \otCroix \hfill \trou[50mm]{No, we weren't.}
\end{questions}
\end{anglais}

\exercice[ecrire,niveau=2]{Notre rue autrefois}{Complétez avec \en{there was}, \en{there were},
\en{there wasn't} ou \en{there weren't}.}
\begin{anglais}
\begin{tcolorbox}[otBase,colback=otNavyTint,colframe=otNavyTint,fontupper=\fontsize{11}{19}\selectfont]
Fifty years ago, \trou[34mm]{there was} a cinema on our street, and \trou[34mm]{there were} two
bakeries. \trou[34mm]{There wasn't} a supermarket, and \trou[34mm]{there weren't} any coffee shops.
But \trou[34mm]{there were} children everywhere, and \trou[34mm]{there was} music in the evening!
\end{tcolorbox}
\end{anglais}

\exercice[lire,niveau=2]{New York, d'hier à aujourd'hui}{Lisez la frise, puis indiquez si chaque phrase
est vraie (V) ou fausse (F).}
\begin{center}\otFriseNY\end{center}
\begin{anglais}
\begin{questions}
\item In 1625, New York was a Dutch town. \hfill\trou[9mm]{V}
\item Its name was New Amsterdam. \hfill\trou[9mm]{V}
\item The Brooklyn Bridge was new in 1931. \hfill\trou[9mm]{F}
\item The Statue of Liberty was a gift from England. \hfill\trou[9mm]{F}
\item Ellis Island was the door to America. \hfill\trou[9mm]{V}
\item The Empire State Building was there in 1886. \hfill\trou[9mm]{F}
\end{questions}
\end{anglais}

\exercice[ecrire,niveau=2]{Corrigez l'erreur}{Chaque phrase contient une erreur. Récrivez-la
correctement.}
\begin{questions}
\item I were at the party. \hfill \enb{\trou[66mm]{I was at the party.}}
\item They was very tired. \hfill \enb{\trou[66mm]{They were very tired.}}
\item Was you at work yesterday? \hfill \enb{\trou[66mm]{Were you at work yesterday?}}
\item I am born in Lyon. \hfill \enb{\trou[66mm]{I was born in Lyon.}}
\item She was here there is two days. \hfill \enb{\trou[66mm]{She was here two days ago.}}
\end{questions}

\exercice[ecouter,niveau=1,piste=17.10]{\en{was} ou \en{were} ?}{Écoutez chaque phrase, dite deux fois.
Cochez la forme que vous entendez.}
\begin{questions}[3]
\item \choix{was}\quad\choix{were}
\item \choix{was}\quad\choix{were}
\item \choix{was}\quad\choix{were}
\item \choix{was}\quad\choix{were}
\item \choix{was}\quad\choix{were}
\item \choix{was}\quad\choix{were}
\end{questions}

\exercice[ecouter,ecrire,niveau=2,piste=17.11]{Le mystère du gâteau}{Le gâteau au chocolat de
Mr.\ Carter a disparu ! Écoutez l'enquête et complétez.}
\begin{anglais}
\begin{questions}
\item This morning, the cake was in the \trou[34mm]{kitchen}.
\item Jack says: « I was at the \trou[34mm]{post office}. »
\item Emma was in the meeting room with a \trou[34mm]{client}.
\item Léa was at her \trou[34mm]{desk}, with Jack.
\item So Jack \trou[34mm]{wasn't} at the post office!
\end{questions}
\end{anglais}

\exercice[ecouter,ecrire,niveau=2,piste=17.12]{Dictée à trous}{Mrs.\ Lopez demande à Léa comment
s'est passé Noël. Écoutez et complétez.}
\begin{anglais}
\noindent\begin{tabular}{@{}l@{\hspace{6mm}}l@{}}
\otNomDit{lopez} & Hi, Léa! How \trou{was} Christmas?\\[1.8mm]
\otNomDit{lea}   & It was wonderful! My parents \trou{were} here for two weeks.\\[1.8mm]
\otNomDit{lopez} & \trou{Were} they happy in New York?\\[1.8mm]
\otNomDit{lea}   & Yes, they \trou{were}. But it was very \trou{cold}!\\[1.8mm]
\otNomDit{lopez} & And \trou{where} were you on New Year's Eve?\\[1.8mm]
\otNomDit{lea}   & In Times Square. There \trou{were} so many people!\\[1.8mm]
\otNomDit{lopez} & Ha! I was there fifty years \trou{ago}!\\
\end{tabular}
\end{anglais}

\exercice[ecrire,niveau=2]{Traduisez}{Traduisez en anglais.}
\begin{questions}[1][french]
\item Où étais-tu hier soir ? — J'étais à la maison.\lignes{1}
\item C'était comment, le concert ? — C'était super !\lignes{1}
\item Ils n'étaient pas en retard.\lignes{1}
\item Tu étais fatigué ? — Oui. \textit{(réponse courte)}\lignes{1}
\item Il y avait beaucoup de monde.\lignes{1}
\item Je suis né à Lyon, il y a trente ans.\lignes{1}
\item Il n'y avait pas de supermarché.\lignes{1}
\item La semaine dernière, il faisait froid.\lignes{1}
\end{questions}

\exercice[ecrire,niveau=3]{Quand j'étais petit}{En cinq phrases, parlez de votre enfance : où vous êtes
né, où était votre maison, ce qu'il y avait dans votre ville, si vous étiez heureux.}
\lignes{5}

\exercice[ecouter,parler,niveau=3,piste=17.13]{À vous de parler !}{Sam vous pose des questions sur le
passé. Répondez \textbf{à voix haute} pendant chaque pause.}
\begin{center}\otTableau
\begin{tabular}{@{}>{\otSemi\fontsize{8.6}{10}\selectfont\color{otNavy}}l@{\hspace{3mm}}>{\enOn}p{70mm}<{\enOff}@{\hspace{6mm}}p{78mm}@{}}
 & \otTH{Sam dit\ldots} & \otTH{Vous\ldots}\\\midrule
1. & Where were you last night? & dites où vous étiez.\\
2. & How was your weekend? & dites comment c'était.\\
3. & Were you at work yesterday? & répondez par une réponse courte.\\
4. & Where were you born? & donnez la ville et l'année.\\
5. & Was there a cinema in your town? & répondez et précisez.\\\bottomrule
\end{tabular}\end{center}

%% ---------------------------------------------------------------------
\rubrique{bilan}{Faites le point}

\begin{essentiel}
\begin{minipage}[t]{0.54\linewidth}
\otEssTitre{Les phrases clés}
\renewcommand{\arraystretch}{1.32}
\begin{tabular}{@{}>{\enOn\otSemi\fontsize{9.4}{11}\selectfont}l<{\enOff}@{\hspace{3mm}}>{\fontsize{9.2}{11}\selectfont\color{white!78!otNavy}}l@{}}
Where were you?            & Où étais-tu ?\\
I was at home.             & J'étais à la maison.\\
It was great!              & C'était super !\\
There were a lot of people. & Il y avait du monde.\\
I was born in Lyon.        & Je suis né(e) à Lyon.\\
two years ago              & il y a deux ans\\
\end{tabular}
\end{minipage}\hfill
\begin{minipage}[t]{0.41\linewidth}
\otEssTitre{\en{was} · \en{were}}
{\fontsize{9.2}{11.6}\selectfont \en{I / he / she / it was} · \en{you / we / they were}\par
\en{wasn't}, \en{weren't} · \en{Was she…? Were you…?}\par}
\otEssTitre{Le passé}
{\fontsize{9.2}{11.6}\selectfont \en{yesterday} · \en{last week} (sans \en{the}) · \en{two days ago}\par}
\otEssTitre{Prononciation}
{\fontsize{9.2}{11.6}\selectfont \en{I was at home} \phon{wəz} · \en{Yes, I was} \phon{wʌz}\par}
\end{minipage}
\end{essentiel}

\begin{jesais}
\item dire où j'étais et comment c'était ;
\item poser une question avec \en{Was…?} ou \en{Were…?} et y répondre ;
\item décrire un lieu autrefois avec \en{there was}, \en{there were} ;
\item situer un événement avec \en{yesterday}, \en{last week}, \en{ago} ;
\item dire où et quand je suis né ;
\item prononcer \en{was} et \en{were} sous leur forme faible.
\end{jesais}

\clearpage
\begin{testexpress}
\qcm[a]{« J'étais à la maison » se dit :}{\en{I was at home.}}{\en{I were at home.}}{\en{I am at home.}}
\qcm[b]{« Où étiez-vous ? » se dit :}{\en{Where was you?}}{\en{Where were you?}}{\en{Where did you be?}}
\qcm[c]{\en{Were they happy?} — Réponse négative :}{\en{No, they wasn't.}}{\en{No, they didn't.}}{\en{No, they weren't.}}
\qcm[a]{« Il y avait deux cinémas » se dit :}{\en{There were two cinemas.}}{\en{There was two cinemas.}}{\en{It had two cinemas.}}
\qcm[b]{« il y a deux ans » se dit :}{\en{there are two years}}{\en{two years ago}}{\en{ago two years}}
\qcm[c]{« Je suis né en 1998 » se dit :}{\en{I am born in 1998.}}{\en{I'm born in 1998.}}{\en{I was born in 1998.}}
\qcm[a]{« la semaine dernière » se dit :}{\en{last week}}{\en{the last week}}{\en{the week last}}
\qcm[c]{Dans \en{I was at home}, \en{was} se prononce :}{\phon{wɑːz}, fort}{\phon{wɪz}}{\phon{wəz}, faible}
\qcm[b]{En 1625, New York s'appelait :}{New London}{New Amsterdam}{New York déjà}
\qcm[a]{Ellis Island était :}{la porte d'entrée des immigrants}{un grand pont}{le premier gratte-ciel}
\end{testexpress}

\noindent{\fontsize{9.6}{12}\selectfont\color{otMuted}Niveau visé (CECRL, niveau A1) : « peut dire où il
était et décrire très simplement une situation passée ». Les corrigés de l'unité sont
page~\pageref{cor-u17-1} ; les textes des pistes d'exercices, page~\pageref{trans-17.10}.\par}
